
Una medida de riesgo es coherente si satisface, entre otras propiedades, la \textbf{subaditividad}:

\[
\rho(X+Y)\leq \rho(X)+\rho(Y)
\]

Esta propiedad implica que la diversificación no debe aumentar el riesgo.

Supóngase que \(X\) y \(Y\) son pérdidas independientes con distribución:

\[
X=
\begin{cases}
0, & \text{con probabilidad } 0.99 \\
10, & \text{con probabilidad } 0.01
\end{cases}
\qquad
Y=
\begin{cases}
0, & \text{con probabilidad } 0.99 \\
10, & \text{con probabilidad } 0.01
\end{cases}
\]

Entonces, al nivel \(99\%\),

\[
VaR_{0.99}(X)=0, \qquad VaR_{0.99}(Y)=0
\]

Ahora, para la suma \(X+Y\):

\[
X+Y=
\begin{cases}
0, & \text{con probabilidad } 0.9801 \\
10, & \text{con probabilidad } 0.0198 \\
20, & \text{con probabilidad } 0.0001
\end{cases}
\]

Como

\[
P(X+Y\leq 0)=0.9801<0.99
\]

y

\[
P(X+Y\leq 10)=0.9999>0.99
\]

se tiene que

\[
VaR_{0.99}(X+Y)=10
\]

Por tanto,

\[
VaR_{0.99}(X+Y)=10 > VaR_{0.99}(X)+VaR_{0.99}(Y)=0
\]

lo cual viola la subaditividad. Así, el \textbf{VaR no es coherente}.

En cambio, para cada posición individual,

\[
CVaR_{0.99}(X)=10, \qquad CVaR_{0.99}(Y)=10
\]

por lo que

\[
CVaR_{0.99}(X)+CVaR_{0.99}(Y)=20
\]

Para la cartera combinada, el peor \(1\%\) de los casos está formado por una probabilidad de \(0.0001\) con pérdida \(20\) y una probabilidad de \(0.0099\) con pérdida \(10\). Entonces,

\[
CVaR_{0.99}(X+Y)=\frac{20(0.0001)+10(0.0099)}{0.01}=10.1
\]

Así,

\[
CVaR_{0.99}(X+Y)=10.1 \leq 20
\]

y en consecuencia,

\[
CVaR_{0.99}(X+Y)\leq CVaR_{0.99}(X)+CVaR_{0.99}(Y)
\]

Por ello, mientras el VaR viola la subaditividad, el \textbf{CVaR sí la satisface} y por tanto \textbf{sí es una medida coherente de riesgo}.